\dykband{EXERCISES IN SWAFL}
{\small Open SWAFL (desktop or online) with the benchmark project. The chapter of the \href{https://github.com/halilunaziru73-creator/Coupled-Water-Fluxes-Pore-Scale-Dynamics-Precision-Irrigation-Mulching-Feedback-Loops/raw/main/Software/v1.1.0/SWAFL_Desktop/docs/manual_v1.1.0/SWAFL_Manual_v1.1.0.pdf}{SWAFL manual} that explains each task is given in brackets; record your values in a table and compare them with what is expected (\emph{Check}).}\par\vspace{3pt}
\begin{exlist}
\item {\color{khblue}[Manual, Chapter 12]} Open \emph{Drip irrigation design}; for the loam and the sand, record the bulb radius and depth for the same emitter.\par\noindent{\small\itshape Check: the sand bulb is narrower and deeper.}
\item {\color{khblue}[Manual, Chapter 12]} Change the lateral length in steps of 20 m and record the emission uniformity. Find the longest acceptable lateral.\par\noindent{\small\itshape Check: the uniformity falls with length.}
\item {\color{khblue}[Manual, Chapter 12]} Use \emph{To size a pump} and record the head and the power for the farm.\par\noindent{\small\itshape Check: the pump head includes friction and pressure.}
\item {\color{khblue}[Manual, Chapter 18]} Compare the three systems on this field and record which one the program prefers and why.\par\noindent{\small\itshape Check: the comparison lists water, energy and uniformity.}
\end{exlist}
